\textbf{7.} Sea $p:E\to B$ una función recubridora, con $E$ conexo por caminos. Demuestre que si $B$ es simplemente conexo entonces $p$ es un homeomorfismo.
\begin{proof}
Supongamos que $B$ es simplemente conexo. Sea $b_0\in B$ cualquiera. Como $p$ es sobreyectiva, existe $e_0\in E$ tal que $p(e_0)=b_0$. Consideremos $\phi:\pi_1(B,b_0)\to p^{-1}(b_0)$, la correspondencia de levantamiento dada por $p$ respecto a $e_0$. Como $E$ es conexo por caminos, $\phi$ es sobreyectiva. Como $B$ es simplemente conexo entonces $|\pi_1(B,b_0)|=1$, luego $\phi$ es trivialmente inyectiva. Así, $\phi$ es una biyección, de donde $|p^{-1}(b_0)|=|\pi_1(B,b_0)|=1$. Como $b_0$ es arbitrario en $B$, de lo anterior se tiene que $p$ es biyectiva. Como además $p$ es continua y abierta por ser una función recubridora, se sigue que $p$ es un homeomorfismo.
\end{proof}\clearpage
